\begin{abstract}
\noindent Sea $Q(n)$ el número de colocaciones de $n$ reinas sin ataques en un tablero de
$n\times n$. Simkin determinó $\log Q(n)$ salvo un error $o(n)$ mediante un problema variacional. Un
trabajo anterior construyó, a partir de una entropía de Bethe, el predictor canónico
$Q_{\rm pred}(n)$ y obtuvo su término constante $C_{\rm pred}$, pero dejó abierto si éste es también
el término constante del conteo exacto. Probamos que lo es:
\[
\begin{aligned}
 \log Q(n)-\log n!-n\Phi &\longrightarrow C_{\rm pred},\\
 C_{\rm pred} &=\tfrac12\log2-\tfrac12-\tfrac32I_2-\tfrac12\log{\det}_2(I+J).
\end{aligned}
\]
Aquí $\Phi=s_1+1$, donde $s_1$ es el valor óptimo en masa uno, que alcanza el perfil de masa uno;
$I_2$ es la integral del cuadrado de ese perfil y $J$ es el operador de Hilbert--Schmidt construido
a partir de él. Equivalentemente, $Q(n)\sim e^{C_{\rm pred}}\,n!\,e^{n\Phi}$ y
$\log Q(n)-\log Q_{\rm pred}(n)\to0$; la cota $O(1/n)$ para esta diferencia, conjeturada en ese
trabajo, sigue abierta.

\smallskip\noindent La prueba compara dos leyes para el proceso que inserta una reina por fila, en
un orden uniforme de las filas: la ley uniforme, que revela una solución uniforme, y la ley
adaptativa, que elige cada casilla según el predictor del tablero resultante. Ese trabajo probó que
el predictor satisface la recursión del conteo salvo un factor $e^{O(k^{-2})}$ en dos clases de
tableros con $k$ filas vacías. Con esa consistencia, una identidad entrópica reduce la comparación
de los conteos al tablero en que se detiene el proceso. Bajo la ley adaptativa, los tableros
satisfacen hasta que quedan $K$ filas vacías unas condiciones que los mantienen en esas clases,
salvo probabilidad $O(e^{-cK})+o(1)$, y acotamos el coste en el primer tablero que deja de
satisfacerlas. Bajo la ley adaptativa, la probabilidad de llegar a $K$ filas vacías sin una salida
y de que dos cruces vacíos compartan diagonal es $O(K^4/n)$. Un cruce vacío es la intersección de
una fila vacía y una columna vacía. Cuando el proceso llega a $K$ sin una salida y ningún par de
cruces vacíos comparte diagonal, el conteo es un permanente, al que aplicamos la estimación de
McCullagh en la forma uniforme probada por Divoux y Wang.

\smallskip\noindent Numéricamente, $C_{\rm pred}\approx0{,}9401$, sin cota de error certificada; los
conteos exactos para ${16\le n\le27}$ son compatibles con el teorema, que no da velocidad de
convergencia.

\medskip\noindent
\textbf{Palabras clave:} problema de las $n$ reinas; enumeración asintótica; término constante;
entropía relativa; permanente.

\noindent
\textbf{MSC 2020:} 05A16; 05A05; 60C05; 94A17; 15A15.
\end{abstract}
